% !TITLE 明月皎夜光
% !AUTHOR 无名氏
% !DYNASTY 两汉
% !GENRE 五言古诗
% !ORDER 5
% !SOURCE 《古诗十九首》；收于《文选》
% !EDITION 《古诗十九首集释》，隋树森，中华书局（2020年版为候选）
% !EDITION_STATUS pending

\begin{poem}
明月皎夜光，促织鸣东壁\textsuperscript{1}。
玉衡指孟冬\textsuperscript{2}，众星何历历\textsuperscript{3}。
白露沾野草，时节忽复易\textsuperscript{4}。
秋蝉鸣树间，玄鸟逝安适\textsuperscript{5}？
昔我同门友\textsuperscript{6}，高举振六翮\textsuperscript{7}。
不念携手好，弃我如遗迹\textsuperscript{8}。
南箕北有斗，牵牛不负轭\textsuperscript{9}。
良无盘石固\textsuperscript{10}，虚名复何益？
\end{poem}

\begin{annotations}
\notetext{1}{促织：蟋蟀。}
\notetext{2}{玉衡：北斗七星中的第五星。古人以斗柄所指的方位确定月份，玉衡指向孟冬方位，即已入初冬。孟冬：冬季的第一个月。}
\notetext{3}{历历：分明可数。}
\notetext{4}{时节忽复易：季节忽然又变换了。易，改变。}
\notetext{5}{玄鸟：燕子。逝安适：飞到哪里去了。适，往。}
\notetext{6}{同门友：同学、同窗。}
\notetext{7}{六翮：鸟翅上的六根强羽，代指翅膀。高举振六翮，比喻仕途得意、青云直上。}
\notetext{8}{遗迹：行人留下的脚印。}
\notetext{9}{南箕北有斗，牵牛不负轭：箕宿像簸箕而不能簸谷，斗宿像斗而不能舀酒，牵牛星叫牛却不拉车。语出《诗经·大东》。轭，驾车时套在牛马颈上的横木。}
\notetext{10}{良：确实。盘石：大石。固：坚固。}
\end{annotations}

\begin{appreciation}
《明月皎夜光》是古诗十九首的第七首，写一个秋夜里被旧友抛弃的人。

前八句全是景。月光皎洁，蟋蟀在东墙下叫；北斗的玉衡星指向孟冬的方位，满天星斗历历分明；白露打湿了野草，季节忽然又换了一轮；秋蝉还在树间鸣叫，燕子却已经不知飞到哪里去了。

这些句子看上去只是记时令，其实每一句都在催人。促织鸣，是秋；玉衡指孟冬，是秋将尽而冬已至；白露沾草，是寒气下来了；蝉鸣将歇，玄鸟南迁，是一年里最后一点生气也在退场。诗人一夜未眠，从月上看到星移，从秋看到冬，时间就在他头顶上一寸一寸地走。

然后笔锋一转：“昔我同门友，高举振六翮。”从前一起读书的朋友，如今高高飞起，翅膀一振就上了天。“六翮”是鸟翅上的六根强羽，用来说人，就是仕途得意、平步青云。可是这位飞黄腾达的同窗，“不念携手好，弃我如遗迹”——不念当年携手的情分，把我像路边的脚印一样丢在身后。

“遗迹”两个字下得极冷。脚印是走路时留在身后的东西，走过就忘了，连回头看都不必。人被旧友这样对待，比被敌人中伤还要难堪，因为你连被记恨的资格都没有，你只是被忘记了。

末四句连用三个星名：“南箕北有斗，牵牛不负轭。”南方的箕宿，形状像簸箕而不能簸谷；北方的斗宿，形状像斗而不能舀酒；天河边的牵牛，名字叫牛却不拉车。这三个比喻出于《诗经·大东》，原诗借它们讽刺周王室有名无实、徒占其位。诗人把它们借过来，说的却是那位“同门友”：名分上是朋友，实际上一点朋友的事也不做。

“良无盘石固，虚名复何益？”——交情本来就没有磐石那样坚固，空有一个朋友的名分，又有什么用处？

《大东》的箕斗牵牛，怨的是天下；这首诗的箕斗牵牛，怨的是一个人。同一种写法，从西周到东汉，把国家大事缩小成了一桩私人的伤心事。可也正因为小，它才扎得深。

天快亮了。星星还在原来的位置上，那个曾经一起读书的人，已经不在了。
\end{appreciation}
